% Generated draft: requires mathematical and source review.
\documentclass[11pt]{article}
\usepackage[margin=25mm]{geometry}
\usepackage{fontspec}
\setmainfont{Latin Modern Roman}
\usepackage{amsmath,amssymb,mathtools,unicode-math}
\setmathfont{Latin Modern Math}
\usepackage{enumitem,needspace,xurl,hyperref}
\setlength{\parindent}{0pt}
\setlength{\parskip}{5pt}
\setlength{\emergencystretch}{3em}
\newcommand{\N}{\mathbb{N}}
\newcommand{\Z}{\mathbb{Z}}
\newcommand{\Q}{\mathbb{Q}}
\newcommand{\R}{\mathbb{R}}
\newcommand{\C}{\mathbb{C}}
\newcommand{\F}{\mathbb{F}}
\newcommand{\A}{\mathbb{A}}
\newcommand{\PP}{\mathbb{P}}
\newcommand{\E}{\mathbb{E}}
\newcommand{\eps}{\varepsilon}
\newcommand{\dd}{\,\mathrm{d}}
\begin{document}
\section*{UnsolvedMath ID: 30000163}
\textbf{Source problem code:} catalogue entry
\subsection{Definitions and mathematical statement}
For each dimension $n\in\{24,8\}$, seek a real nonzero admissible function $f:\R^n\to\R$, with Fourier transform convention
\[
 \widehat f(t)=\int_{\R^n}f(x)e^{-2\pi\sqrt{-1}\,x\cdot t}\,\dd x,
\]
satisfying the exact displayed source conditions
\[
 f(0)=\widehat f(0)>0,\qquad f(x)\le0\ (|x|\ge2),\qquad
 \widehat f(t)\ge0\ (t\in\R^n).
\]
There is a separate function in each dimension. Admissibility means the analytic hypotheses needed for the report's Fourier linear-programming bound; real Schwartz functions are a standard sufficient class, but the extract does not give the exact broader admissibility definition. Euclidean norm is used.

The stated purpose is to prove the sharp sphere-packing bounds for the Leech lattice and $E_8$. The normalization must be checked. With the displayed Fourier convention, the standard bound at exclusion radius $r$ is
\[
 \delta\le\operatorname{vol}(B_n(r/2))\frac{f(0)}{\widehat f(0)}.
\]
Thus $r=2$ and ratio one give $\delta\le\operatorname{vol}(B_n(1))$. This agrees with the usual unimodular Leech normalization, whose shortest nonzero vector has length two. In dimension eight, the usual unimodular $E_8$ has shortest-vector length $\sqrt2$. Its sharp density is $2^{-4}\operatorname{vol}(B_8(1))$. Consequently sharpness in dimension eight uses either $r=\sqrt2$ with ratio one, or $r=2$ with ratio $2^{-4}$ under the standard convention. The literal $r=2$, ratio-one conditions in the supplied clean statement do not by themselves yield that sharp $E_8$ density. The source parameters are preserved rather than silently changed; the intended Theorem 4 convention needs reconciliation.

The literature assessment checked on 21 August 2026 cites only the original report and says no later resolution was located. That is the supplied search outcome, not independent verification of current openness. It must not override the normalization issue or be promoted to an authoritative present-status conclusion.
\subsection{Short English statement}
Construct the specified Fourier auxiliary functions in dimensions twenty-four and eight. The supplied parameters match the Leech scaling, but the claimed sharp $E_8$ consequence needs a radius or Fourier-ratio normalization correction; the original parameters remain visible for review.
\subsection{Sources}
\begin{itemize}
\item \emph{Gitter und Anwendungen} (2005), Oberwolfach Reports, DOI 10.4171/owr/2005/01, pp. 5--67, \url{https://doi.org/10.4171/owr/2005/01}. The original question refers to its Theorem 4.
\item The sole source in the literature assessment checked 21 August 2026 is the original OWR report, \url{https://doi.org/10.4171/owr/2005/01}. No later result is supplied in the record.
\item Local UnsolvedMath registry, permanent ID $30000163$, exact source code \texttt{catalogue entry}, created 25 August 2026 and updated 16 September 2026. Its clean statement is labelled corrected and verified, with Fourier hypotheses, dimensions and radius restored; its open label remains supplied dated metadata.
\end{itemize}
\end{document}
